\chapter{Background}
\label{chap:background}

% ─────────────────────────────────────────────────────────────────────────────
\section{Concepts Overview}
\label{sec:concepts}
% ─────────────────────────────────────────────────────────────────────────────

This section provides an overview of the key concepts and terminology necessary to
understand the design and motivation of a multi-agent system for AI clinical decisioning
via automation workflows.  The concepts introduced here form the foundation for the
literature review that follows, covering the limitations of single-model AI in
healthcare, the principles of multi-agent
system design, the role of large language models in clinical
applications, and the fundamentals of automation workflow
orchestration.

% ─────────────────────────────────────────────────────────────────────────────
\subsection{The Limitations of Single-Model AI in Clinical Settings}
\label{subsec:single_model_limits}
% ─────────────────────────────────────────────────────────────────────────────

The following subsection examines the three core limitations of single-model AI in
clinical settings --- the knowledge gap problem, the single-pass reasoning constraint,
and cognitive reasoning failures --- that collectively motivate the transition to
multi-agent architectures.

Artificial Intelligence has increasingly been applied to medical tasks, including
diagnosis support, drug recommendation, and clinical documentation.  Early approaches
relied on single large language models (LLMs) --- systems trained on vast corpora of
text that can generate clinically plausible responses to natural language
prompts~\cite{Abbasian2024openCHA}.  While these models have demonstrated impressive
performance on structured question-answering tasks, their deployment in real-world
clinical settings exposes a set of persistent limitations~\cite{Yan2024ClinicalLab}.

The most significant limitation is the \textit{knowledge gap problem}.  A single LLM,
regardless of scale, encodes a fixed set of parametric knowledge from its training data
and has no mechanism to access real-time patient information, query external clinical
databases, or perform analytical computations.  Tang~et~al.~\cite{Tang2024MedAgents}
demonstrated empirically that 77\% of errors made by LLMs on medical tasks are
attributable to domain knowledge gaps rather than failures in reasoning --- a finding
that directly motivates the multi-agent approach, in which specialist agents pool
complementary knowledge to compensate for any individual agent's deficiencies.

A second limitation is the \textit{single-pass reasoning constraint}.  Standard LLMs
produce responses in a single forward pass, which is fundamentally unsuitable for
clinical decisioning tasks that require iterative information gathering, differential
diagnosis across multiple hypotheses, and adaptation to evolving patient presentations.
Shang~et~al.~\cite{Shang2025DynamiCare} showed that dynamic, multi-round clinical
decision-making significantly outperformed single-pass approaches on complex patient
cases, establishing iterative reasoning as a structural requirement for clinical AI.

Third, single models are prone to \textit{cognitive reasoning failures}, including
anchoring bias --- the tendency to over-weight an initial diagnosis ---
confirmation bias, and availability bias.  Ke~et~al.~\cite{Ke2024CognitiveBias}
demonstrated through a controlled simulation study that these biases occur
systematically in single LLM-based clinical agents and are substantially reduced when a
multi-agent framework with a dedicated critical evaluator is used instead.

% ─────────────────────────────────────────────────────────────────────────────
\subsection{Multi-Agent Systems: Principles and Architecture}
\label{subsec:mas_principles}
% ─────────────────────────────────────────────────────────────────────────────

Having identified the limitations of single-model approaches, this subsection introduces
the multi-agent system paradigm as an architectural response, covering the core design
dimensions of topology, coordination, and agent
specialisation.

A multi-agent system (MAS) is a computational framework in which multiple autonomous
agents --- each capable of perceiving its environment, reasoning over inputs, and taking
actions --- collaborate to accomplish tasks that would be too complex or too broad for
any single agent to handle alone~\cite{Kim2024MDAgents}.  In the context of LLM-based
clinical AI, each agent is a language model instance assigned a specific role, knowledge
domain, or workflow function.  Agents communicate through structured message-passing
protocols and may have access to distinct tool sets, knowledge bases, or data
sources~\cite{Abbasian2024openCHA}.

The architectural design of a MAS can take several forms depending on how agents are
coordinated.  In a \textbf{centralised} architecture, all agents report to a single
orchestrator agent that plans, delegates, and synthesises outputs.  In a
\textbf{decentralised} architecture, agents communicate peer-to-peer without a central
coordinator.  \textbf{Hierarchical} architectures combine elements of both, grouping
agents into tiers with escalation pathways between
levels~\cite{Kim2025TAO}.  Chen~et~al.~\cite{Chen2025MedSentry} demonstrated empirically
that the choice of topology has a direct and measurable effect on clinical safety --- a
centralised MAS was found to be the most vulnerable to adversarial manipulation, while a
decentralised topology was the most resilient.

Agent specialisation is a further architectural dimension.  Rather than deploying
identical generalist agents, a specialised MAS assigns distinct roles --- such as a
Retrieval Agent, a Reasoning Agent, a Synthesis Agent, or a Critical Evaluator --- to
separate agents, enabling each to be optimised for its specific function within the
workflow~\cite{Zhou2024ZODIAC,Yan2024ClinicalLab}.  This specialisation principle is
operationalised across nearly all frameworks reviewed in this chapter.

% ─────────────────────────────────────────────────────────────────────────────
\subsection{Large Language Models as Clinical Agents}
\label{subsec:llm_agents}
% ─────────────────────────────────────────────────────────────────────────────

This subsection describes how large language models serve as the reasoning backbone of
individual agents within a multi-agent clinical system, and introduces
retrieval-augmented generation as the primary mechanism for extending agent knowledge
beyond parametric training data.

Large language models are neural network architectures trained on large-scale text
corpora using self-supervised objectives, enabling them to generate coherent,
contextually appropriate text across a wide range of domains~\cite{Tang2024MedAgents}.
In the context of multi-agent clinical systems, LLMs function as the reasoning backbone
of individual agents: given a structured prompt describing the agent's role, the clinical
context, and any inputs from other agents, the LLM generates the agent's response or
action~\cite{Kim2024MDAgents}.

Contemporary LLMs --- including GPT-4, Claude, Gemini, and open-source models such as
Llama-3 --- differ in their parameter counts, training data compositions,
instruction-following capabilities, and factual accuracy on medical benchmarks.
Chen~et~al.~\cite{Chen2024RareAgents} demonstrated that a Llama-3.1-70B-backed
multi-agent system could outperform GPT-4o on rare disease diagnosis tasks, establishing
that agent architecture and coordination quality can compensate for differences in
individual model capability.

The use of \textit{retrieval-augmented generation} (RAG) extends the effective knowledge
of an LLM agent beyond its parametric training data by retrieving relevant documents,
clinical guidelines, or patient records at inference time.
Choi~et~al.~\cite{Choi2024MALADE} demonstrated that RAG integration within a
multi-agent pharmacovigilance system reduced hallucinated drug--event associations by
64\% compared to non-RAG baselines --- a clinically critical improvement that has been
replicated across multiple reviewed frameworks~\cite{Lu2024TriageAgent,Chen2024RareAgents}.

% ─────────────────────────────────────────────────────────────────────────────
\subsection{Automation Workflow Design in Clinical AI}
\label{subsec:workflow_design}
% ─────────────────────────────────────────────────────────────────────────────

This subsection introduces the concept of automation workflow design --- the systematic
decomposition of clinical tasks into structured agent pipelines --- and identifies the
four inter-related design challenges that any well-designed clinical workflow must
address.

Automation workflow design refers to the systematic decomposition of a complex clinical
task into a structured sequence of sub-tasks, each assigned to a specialised agent or
tool, with defined handoff protocols and quality-control mechanisms between
stages~\cite{Barra2025Agentic}.  The workflow pattern is the core architectural
contribution that distinguishes different MAS frameworks from one another.

The foundational workflow pattern --- gather, analyse, synthesise --- was established by
Tang~et~al.~\cite{Tang2024MedAgents} in MedAgents and has been extended by subsequent
frameworks in various directions: adding parallelisation of independent sub-tasks
\cite{Barra2025Agentic}, introducing iterative multi-round loops with dynamic stopping
criteria~\cite{Lu2024TriageAgent}, adding error-recovery sub-loops~\cite{Wu2024EHRFlow},
and extending to dynamic multi-round patient
interaction~\cite{Shang2025DynamiCare}.  Bedi~et~al.~\cite{Bedi2025OptParadox}
provided the critical theoretical insight that workflow coordination quality ---
specifically the compatibility of agents across handoff boundaries --- is a stronger
determinant of overall system performance than the quality of individual components.

A well-designed clinical automation workflow must therefore address four inter-related
design challenges: task decomposition into appropriately scoped agent roles; coordination
protocols governing inter-agent communication; conflict resolution mechanisms for
handling agent disagreements; and safety mechanisms ensuring outputs remain within
clinically acceptable bounds~\cite{Fan2024AIHospital,Wang2025Catfish,Chen2025MedSentry}.
The following literature review examines how the current body of research has addressed
each of these challenges.

% ─────────────────────────────────────────────────────────────────────────────
\section{Search Strategy}
\label{sec:search}
% ─────────────────────────────────────────────────────────────────────────────

A structured search of the academic literature was conducted across five databases:
arXiv, PubMed, ACL Anthology, IEEE Xplore, and Semantic Scholar.  Search terms were
drawn from the core concepts of the thesis: ``multi-agent system AND clinical decision'',
``LLM agent AND medical reasoning'', ``automation workflow AND clinical AI'', ``multi-agent
AND electronic health records'', and ``agent collaboration AND healthcare''.  Results were
filtered to publications from 2024 onwards to ensure relevance to the current state of
the field, with one exception for foundational conceptual papers.

Articles were included if they: (1)~proposed or evaluated a multi-agent or agentic LLM
framework; (2)~applied the framework to a clinical or medical task; (3)~reported
quantitative performance metrics or clinician-validated outcomes; and (4)~were published
in a peer-reviewed venue or as a preprint on arXiv with full methodological disclosure.
Articles were excluded if they described single-model AI systems without agent
coordination, focused exclusively on medical image analysis without a clinical decisioning
component, or were review articles without original empirical contributions.  Following
this process, 22 papers were selected and organised into the five thematic categories
described in the sections that follow.  Table~\ref{tab:inclusion} summarises the
inclusion and exclusion criteria applied.

\begin{table}[H]
\centering
\caption{Inclusion and Exclusion Criteria for Selected Articles}
\label{tab:inclusion}
\begin{adjustbox}{max width=\textwidth}
\begin{tabular}{p{3.0cm} p{5.5cm} p{5.5cm}}
\toprule
\textbf{Criteria} & \textbf{Inclusion} & \textbf{Exclusion} \\
\midrule
Publication Period
  & 2024--2025 (2023 for foundational works)
  & Articles published before 2023 \\[4pt]
System Type
  & Multi-agent or agentic LLM frameworks for clinical/medical tasks
  & Single-model AI systems without agent coordination \\[4pt]
Task Domain
  & Clinical decisioning, medical diagnosis, EHR analysis, clinical workflow automation
  & Medical image analysis without clinical decisioning component \\[4pt]
Evidence Type
  & Quantitative metrics or clinician-validated outcomes with full methodological disclosure
  & Review articles or opinion papers without original empirical contributions \\[4pt]
Language
  & English-language publications
  & Non-English publications \\
\bottomrule
\end{tabular}
\end{adjustbox}
\end{table}
